\section{Imaging, navigation, integrity, and target selection}
\label{sec:navigation-maths}

This section audits proposal Sections 10--12 and 14 and Appendices
A--D, G, H, and L. Algebraic identities are distinguished from statistical
claims, which require the stated noise model. Numerical checks below use
synthetic data, not flight measurements. They verify equations and
implementations under their assumptions; they do not establish the proposed
flight accuracy, reliability, computation time, or terrain-match success.

\subsection{Frames, signs, units, and camera model}

Use a local Cartesian world frame with upward vertical coordinate $z$,
ground plane $z=0$, and gravity $\boldsymbol g=(0,0,-g_0)^{\mathsf T}$.
Let $R_{wb}\in SO(3)$ take body coordinates to world coordinates,
$R_{cb}\in SO(3)$ take body coordinates to camera coordinates, and
$\boldsymbol\ell_b$ denote the camera centre relative to the IMU origin in
body coordinates. The navigation position $\boldsymbol p$ is the IMU
origin. Camera depth is positive in front of the camera. Thus
\begin{align}
 \boldsymbol C&=\boldsymbol p+R_{wb}\boldsymbol\ell_b,
 &R_{cw}&=R_{cb}R_{wb}^{\mathsf T},\\
 \boldsymbol q_c&=R_{cb}\{R_{wb}^{\mathsf T}(\boldsymbol m-\boldsymbol p)
                         -\boldsymbol\ell_b\}
                 =R_{cw}(\boldsymbol m-\boldsymbol C),\\
 \boldsymbol z&=\pi(K\boldsymbol q_c),
 &\pi([a,b,c]^{\mathsf T})&=[a/c,b/c]^{\mathsf T}.
 \label{eq:full-camera}
\end{align}
Here $K$ has positive focal lengths $f_x,f_y$ in pixels, principal point
$(c_x,c_y)$, and, unless stated otherwise, zero skew. Image coordinates
must be undistorted before using this pinhole model, or distortion must be
included in the observation function and its derivative. All attitude
increments and angular rates in calculations are in radians.

The proposal uses $R$ for a world-to-camera rotation in Section 10 and
a body-to-world rotation in Appendix H. Both equations can be valid, but
the same symbol must not be passed between them without taking the
transpose and applying the camera extrinsic rotation. Equation
\eqref{eq:full-camera} removes that ambiguity. A camera lever arm may be
neglected only after its magnitude and associated rotational displacement
have been shown small relative to the error allocation.

\subsection{Ground-plane rectification and field of view}

Write $R_{cw}=[\boldsymbol r_1\ \boldsymbol r_2\ \boldsymbol r_3]$.
For $\boldsymbol m=(X,Y,0)^{\mathsf T}$,
\begin{equation}
 K R_{cw}(\boldsymbol m-\boldsymbol C)
 =\underbrace{K[\boldsymbol r_1\ \boldsymbol r_2\ -R_{cw}\boldsymbol C]}_H
   [X,Y,1]^{\mathsf T}.
 \label{eq:homography-proof}
\end{equation}
Homogeneous scaling does not change the pixel after division by its last
coordinate. This proves the proposed homography directly. Rectification
is inverse image sampling: for each desired ground-grid point, compute its
image pixel through $H$, then interpolate the recorded image there. It
does not require numerically inverting $H$ for every ground pixel.

The map-to-image homography is invertible when the camera is not in the
plane and $K$ is nonsingular. Indeed, because rotations preserve cross
products,
\begin{equation}
 \det H=\det K\,(\boldsymbol r_1\times\boldsymbol r_2)^{\mathsf T}
                   (-R_{cw}\boldsymbol C)
        =-C_z\det K.
\end{equation}
Invertibility alone does not guarantee visibility. Require positive
camera depth and pixels inside the valid image, and exclude rays nearly
parallel to the ground, where division is ill-conditioned. A tilted
camera generally produces a quadrilateral footprint, not the nadir
rectangle. Nonplanar terrain produces parallax that a single homography
cannot remove.

For a nadir camera with \emph{horizontal} field of view $90^\circ$, image
width $1280$, and square pixels, $f_x=f_y=640$ pixels. At height $h$ the
ground width and height are
\begin{equation}
 W=\frac{1280h}{640}=2h,\qquad
 H_g=\frac{800h}{640}=1.25h,\qquad
 r_g=\frac h{640}\ \mathrm{m/pixel}.
\end{equation}
At $h=300\,\mathrm m$ these give $600\times375\,\mathrm m$ and
$0.46875\,\mathrm{m/pixel}$. The vertical field of view is
$2\arctan(400/640)=64.0108^\circ$. These numbers require a horizontal
$90^\circ$ specification; a diagonal or vertical $90^\circ$ lens
specification gives different focal lengths and footprints.

For a camera height error caused by neglecting terrain elevation
$\delta z$, an oblique ray at angle $\phi$ has horizontal intersection
error $\delta z\tan\phi$ to first order. This is geometry, not proof that
terrain errors are independent, Gaussian, or of fixed magnitude across
landmarks.

\subsection{Timing, exposure blur, and tilt--translation separation}

\subsubsection{Timestamp error and exposure}

An image observed at $t$ but assigned attitude at $t+\delta\tau$ has
first-order attitude error $\boldsymbol\omega\delta\tau$. For a central
nadir ray and a pitch or roll rate $\omega$, the corresponding ground
error is $h\omega\delta\tau$. Allocating $b_t$ metres therefore gives
\begin{equation}
 |\delta\tau|\leq\frac{b_t}{h|\omega|}.
\end{equation}
For $b_t=1\,\mathrm m$, $h=300\,\mathrm m$, and
$|\omega|=40\pi/180=0.698132\,\mathrm{rad/s}$, the limit is
$4.77465\,\mathrm{ms}$. The proposal's $5\,\mathrm{ms}$ is a rounded
value, not an exact upper limit. A $0.5\,\mathrm{ms}$ offset contributes
$0.10472\,\mathrm m$ in this central-ray example.

For tilt $\beta$, the nadir intersection is $h\tan\beta$ and its derivative
is $h\sec^2\beta\,\omega$. Motion of the camera centre adds a term
$\boldsymbol v\delta\tau$. In the full pixel model the timing derivative
is $\dot{\boldsymbol z}=H_p\boldsymbol v+H_\theta\boldsymbol\omega$,
with consistent frames and any lever-arm motion included. Thus the
central-ray formula is a design approximation, not a bound for all
pixels and attitudes. Treat unknown timing offset as a shared nuisance
parameter if it is estimated from the same images.

Exposure blur is a different effect. A pitch perturbation at normalized
horizontal coordinate $x$ produces speed
$f(1+x^2)|\omega|$ pixels/s; a yaw rate produces speed
$r_{\mathrm{px}}|\omega_z|$ at image radius $r_{\mathrm{px}}$.
With $f=640$ and $40^\circ/\mathrm s$, a $1\,\mathrm{ms}$ exposure gives
$0.44680$ pixel at the centre and $0.89361$ pixel at the horizontal
$45^\circ$ edge. A $1$ revolution/s yaw rate gives
$2\pi\sqrt{640^2+400^2}(0.001)=4.74204$ pixels at a corner.
Halving exposure halves these first-order values. Finite angle, mixed
rotation, lens distortion, and translation can change the result; actual
exposure feasibility requires measurements of light and sensor noise.

\subsubsection{The one-dimensional separation proof}

For one horizontal section of a nadir image, write
$x_i=X_i/h=\tan\alpha_i$ and $u_i=f x_i$. A camera translation
$\delta p_x$ gives $\delta u_i=-f\delta p_x/h$. A small pitch
$\delta\theta$ gives
\begin{equation}
 \delta u_i=-f(1+x_i^2)\delta\theta.
\end{equation}
Choosing $\delta p_x=h\delta\theta$ matches the central shift and leaves
$-f x_i^2\delta\theta=-f\tan^2\alpha_i\,\delta\theta$, as claimed in
Appendix C.

The useful statistical statement follows by explicitly eliminating
unknown translation. With independent pixel noise of variance
$\sigma_{\mathrm{px}}^2$, the Jacobian row for
$(\delta p_x,\delta\theta)$ is
$[-f/h,-f(1+x_i^2)]$. The Fisher information is its weighted outer-product
sum. Its Schur complement after marginalizing translation is
\begin{align}
 J_{\theta\mid\mathrm{unknown}\ p_x}
 &=\frac{f^2}{\sigma_{\mathrm{px}}^2}
   \left\{\sum_i(1+x_i^2)^2-
                    \frac{[\sum_i(1+x_i^2)]^2}{n}\right\}\\
 &=\frac{f^2}{\sigma_{\mathrm{px}}^2}
                  \sum_i(x_i^2-\overline{x^2})^2.
 \label{eq:tilt-information}
\end{align}
This is positive exactly when the squared off-axis coordinates are not
all equal. If all landmarks are at $x_i=\pm a$, even at the image edges,
the one-dimensional tilt and translation columns are proportional and
the information is zero. A range of off-axis distances is necessary in
this reduced model. If translation is known rather than marginalized,
the information is instead
$f^2\sum_i(1+x_i^2)^2/\sigma_{\mathrm{px}}^2$.

\subsubsection{Full two-dimensional geometry and observability}

For $\boldsymbol q_c=Z[x,y,1]^{\mathsf T}$, with local camera attitude
increment, the normalized-image rotational derivative is
\begin{equation}
 \frac{\partial(x,y)}{\partial\delta\boldsymbol\theta_c}
 =\begin{bmatrix}
 xy&-(1+x^2)&y\\
 1+y^2&-xy&-x
 \end{bmatrix}.
 \label{eq:full-image-rotation}
\end{equation}
Multiplying rows by $f_x,f_y$ gives the pixel derivative. The normalized
translation derivative in camera coordinates is
\begin{equation}
 \frac1Z\begin{bmatrix}-1&0&x\\0&-1&y\end{bmatrix}.
\end{equation}
These follow by differentiating $X/Z,Y/Z$, so the signs are tied to
\eqref{eq:full-camera}. The $xy$ cross terms and the second image
coordinate contain additional information. Appendix C's residual is
therefore the separating information in its one-dimensional example,
not the only information in the full camera model.

For an actual set of landmarks, form the full weighted pose Jacobian.
Partition translation and other nuisance parameters from attitude,
then inspect the Schur complement and its eigenvalues. Wide angular
coverage can improve conditioning, but cannot guarantee a desired
attitude accuracy regardless of feature geometry, unknown height,
extrinsic calibration, terrain relief, or measurement noise. A known
metric planar map and nondegenerate correspondences can locally constrain
pose; finite planar pose ambiguities, association ambiguities, and
weak conditioning still require handling. Velocity and IMU biases require
information over time, not just a single image.

Specific force is not a gravity direction during general acceleration.
The identity $\boldsymbol f_b=R_{wb}^{\mathsf T}
(\ddot{\boldsymbol p}-\boldsymbol g)$ shows that without a model or
additional measurements, arbitrary acceleration can compensate a change
in tilt. For the ideal aligned pendulum, line tension is the only
nongravitational force and its measured direction is fixed in body
coordinates: an accelerometer alone then supplies no instantaneous swing
angle. Swing averaging becomes a useful initialization only under
additional dynamical assumptions, not an observability theorem.

\subsection{Landmark correlation, peak uncertainty, and consensus}

\subsubsection{Normalized cross-correlation and its fast evaluation}

For a template containing $M$ pixels define
$T_0(\boldsymbol x)=T(\boldsymbol x)-\overline T$,
$E_T=\sum T_0^2$, and, at shift $\boldsymbol d$,
\begin{align}
 S_I(\boldsymbol d)&=\sum_{\boldsymbol x}I(\boldsymbol x+\boldsymbol d),
 &S_{I^2}(\boldsymbol d)&=\sum_{\boldsymbol x}I^2(\boldsymbol x+\boldsymbol d),\\
 E_I(\boldsymbol d)&=S_{I^2}(\boldsymbol d)-S_I(\boldsymbol d)^2/M,\\
 C(\boldsymbol d)&=\sum_{\boldsymbol x}T_0(\boldsymbol x)
                                      I(\boldsymbol x+\boldsymbol d).
\end{align}
Because $\sum T_0=0$, subtracting the local image mean in the numerator
does not change $C$. Hence
\begin{equation}
 \gamma(\boldsymbol d)=\frac{C(\boldsymbol d)}{
                         \sqrt{E_T E_I(\boldsymbol d)}}.
 \label{eq:ncc-fast}
\end{equation}
Cauchy--Schwarz proves $|\gamma|\leq1$ when both energies are positive.
Constant patches have zero energy and must be rejected or specially
handled; assigning them a high-confidence match is undefined.

The correlation theorem gives
\begin{equation}
 C=\mathcal F^{-1}\{\overline{\mathcal F(T_0)}\,\mathcal F(I)\}.
\end{equation}
The arrays must be padded for linear correlation, at least to the sum
of image and template dimensions minus one on each axis, unless an
equivalent overlap-save scheme is used. Otherwise wrap-around creates
spurious matches. Crop to valid overlap with a documented lag convention.
Integral images of $I$ and $I^2$ give each rectangular sum with four
lookups after linear-time preparation. Thus the padded numerator costs
$O(N\log N)$ for $N$ padded pixels, while normalization costs $O(N)$;
this does not imply FFTs always beat direct correlation for small
templates. Masks, missing pixels, or shift-dependent template support
require corresponding weighted sums rather than the fixed-support
formula. The fast normalization construction is described in the primary
paper by Lewis.\footnote{J.~P. Lewis, \emph{Fast Normalized
Cross-Correlation}: \url{https://www.ncorr.com/download/publications/lewisfast.pdf}.}

\subsubsection{Subpixel interpolation}

Fit $q(t)=a t^2+b t+c$ through samples $\gamma_{-1},\gamma_0,
\gamma_{+1}$ at $t=-1,0,1$. Solving the three equations gives
\begin{equation}
 a=\frac{\gamma_{-1}-2\gamma_0+\gamma_{+1}}2,
 \qquad b=\frac{\gamma_{+1}-\gamma_{-1}}2,
 \qquad c=\gamma_0.
\end{equation}
The stationary point $-b/(2a)$ is therefore
\begin{equation}
 \delta=\frac{\gamma_{-1}-\gamma_{+1}}
          {2(\gamma_{-1}-2\gamma_0+\gamma_{+1})}.
\end{equation}
Accept it only for a concave, sufficiently curved, isolated interior
maximum and a plausible subpixel offset. A flat or convex fit does not
yield a useful peak location.

Independent axis fits neglect mixed curvature. For a general local fit
$\gamma(\boldsymbol d)=c+\boldsymbol b^{\mathsf T}\boldsymbol d
-\tfrac12\boldsymbol d^{\mathsf T}A\boldsymbol d$ with $A\succ0$,
the joint maximum is $A^{-1}\boldsymbol b$. For example,
$\gamma=1-(\boldsymbol d-\boldsymbol d_*)^{\mathsf T}
Q(\boldsymbol d-\boldsymbol d_*)$ with
\begin{equation}
 Q=\begin{bmatrix}3&1\\1&2\end{bmatrix},\qquad
 \boldsymbol d_*=[0.23,-0.28]^{\mathsf T},
\end{equation}
has separate-axis offsets $(0.136667,-0.165)$ but joint optimum
$(0.23,-0.28)$. This is an exact quadratic counterexample to treating
axis fits as universally exact two-dimensional interpolation.

\subsubsection{What peak curvature does and does not establish}

A dimensionless correlation curvature is not automatically an inverse
covariance. Under a specified image model
\begin{equation}
 \boldsymbol y=a\boldsymbol T(\boldsymbol d)+b\boldsymbol 1
                  +\boldsymbol\varepsilon,
 \qquad\boldsymbol\varepsilon\sim N(0,\sigma_I^2 I),
\end{equation}
let $D=\partial\boldsymbol T/\partial\boldsymbol d$ and
$B=[\boldsymbol T,\boldsymbol1]$. Marginalizing the unknown photometric
gain and offset by a Schur complement gives
\begin{equation}
 J_d=\frac{a^2}{\sigma_I^2}D^{\mathsf T}
 \left[I-B(B^{\mathsf T}B)^{-1}B^{\mathsf T}\right]D.
 \label{eq:photometric-fisher}
\end{equation}
This follows from the Gaussian likelihood Fisher matrix
$\sigma_I^{-2}[aD,B]^{\mathsf T}[aD,B]$. Inverse information is meaningful
only on identifiable directions. If texture varies along one direction,
the other may be unobservable, the usual aperture problem.

The proposal's $c(-\nabla^2\gamma)^{-1}$ is a calibration model, not
an identity. It needs positive definite negative curvature, compatible
coordinate units, and validation against known shifts with the same
texture, blur, resampling, and photometric mismatch. Adjacent scores,
map pixels, and resampled image pixels are correlated. Wrong-peak errors
are association failures and cannot be represented merely by a Gaussian
covariance about the selected peak. A sharp aliased peak may be wrong.
Measurements produced on a pose-rectified image should be mapped back to
original image coordinates with the warp and covariance accounted for,
or use an observation model that explicitly represents their dependence
on the predicted pose.

\subsubsection{RANSAC iteration count and model choice}

If a trial has probability $q$ of selecting a valid all-inlier minimal
set and independent trials are drawn, the probability all $N$ trials
fail is $(1-q)^N$. Requiring success probability at least $p$ gives
\begin{equation}
 N\geq\left\lceil\frac{\log(1-p)}{\log(1-q)}\right\rceil.
 \label{eq:ransac-ceil}
\end{equation}
With independent inlier labels and sample size $k$, $q=w^k$.
With $n$ fixed correspondences, $n_I$ known inliers, and $k$ distinct
points sampled uniformly within each independently redrawn trial,
\begin{equation}
 q=\frac{\binom{n_I}{k}}{\binom n k}.
\end{equation}
Degenerate all-inlier samples lower $q$ further. These formulas are
sampling probabilities, not probabilities that the finally chosen model
is correct. Structured aliases and unmodelled distortions can generate
large wrong consensuses. The original RANSAC paper establishes the
random-consensus framework.\footnote{Fischler and Bolles (1981):
\url{https://doi.org/10.1145/358669.358692}.}

An orientation-preserving two-dimensional similarity obeys
$m_i=z u_i+t$ in complex coordinates. Two distinct $u_i$ determine
\begin{equation}
 z=\frac{m_2-m_1}{u_2-u_1},\qquad t=m_1-z u_1.
\end{equation}
Require distinct source points, nonzero admissible scale, and adequate
baseline relative to noise. For $w=0.5$, $k=2$, and $p=0.999$,
$N\geq\lceil24.011768\rceil=25$. With $24$ trials the success
probability is $1-0.75^{24}=0.998996609$, slightly short of the target.
Using $n=50,n_I=25$ also requires $25$ trials under the distinct-point
formula. A general planar homography needs four point pairs in general
position; the same independent-label assumptions give $108$ trials
for $k=4$.

The residual between two differently rectified plane images is in
general a homography, not a similarity. Similarity consensus is an
approximation justified only when the remaining perspective error is
small enough for its threshold. Alternatively use a homography or a
physical camera-pose hypothesis and its appropriate minimal solver.
Reject geometric degeneracy explicitly and refine the selected model
on its accepted correspondences.

\subsection{Error-state navigation filter with a coherent convention}

The local/right attitude convention used here agrees with the proposal's
Appendix H and the convention analyzed by Sol\`a.\footnote{Joan Sol\`a,
\emph{Quaternion kinematics for the error-state Kalman filter}:
\url{https://arxiv.org/abs/1711.02508}.}
Define $[\boldsymbol a]_\times\boldsymbol b=\boldsymbol a\times
\boldsymbol b$, true-minus-nominal Euclidean errors, and
\begin{equation}
 R_{wb}=\widehat R_{wb}\operatorname{Exp}
                      ([\delta\boldsymbol\theta]_\times).
 \label{eq:right-attitude-error}
\end{equation}
For $\vartheta=\|\boldsymbol\theta\|$, Rodrigues' formula gives
\begin{equation}
 \operatorname{Exp}([\boldsymbol\theta]_\times)
 =I+\frac{\sin\vartheta}{\vartheta}[\boldsymbol\theta]_\times
  +\frac{1-\cos\vartheta}{\vartheta^2}[\boldsymbol\theta]_\times^2,
\end{equation}
with continuous small-angle limits. If scalar-first Hamilton quaternions
are used with $R(q_1\otimes q_2)=R(q_1)R(q_2)$, the increment is
$q_\delta=[\cos(\vartheta/2),
\boldsymbol\theta\sin(\vartheta/2)/\vartheta]$ and is multiplied on the
right. A different quaternion convention requires corresponding changes.

\subsubsection{Propagation and its linearization}

Let the measured specific force and angular rate be
\begin{equation}
 \boldsymbol a_m=\boldsymbol f_b+\boldsymbol b_a+\boldsymbol n_a,
 \qquad\boldsymbol\omega_m=\boldsymbol\omega_b+
                  \boldsymbol b_g+\boldsymbol n_g.
\end{equation}
With $\widehat{\boldsymbol f}=\boldsymbol a_m-\widehat{\boldsymbol b}_a$
and $\widehat{\boldsymbol\omega}=\boldsymbol\omega_m-
\widehat{\boldsymbol b}_g$, nominal propagation is
\begin{equation}
 \dot{\widehat{\boldsymbol p}}=\widehat{\boldsymbol v},\quad
 \dot{\widehat{\boldsymbol v}}=\widehat R_{wb}\widehat{\boldsymbol f}
                       +\boldsymbol g,\quad
 \dot{\widehat R}_{wb}=\widehat R_{wb}[\widehat{\boldsymbol\omega}]_\times.
\end{equation}
Assume biases are random walks with independent zero-mean driving noises
$\boldsymbol n_{ba},\boldsymbol n_{bg}$. Since
$[\delta\boldsymbol\theta]_\times\widehat{\boldsymbol f}
=-[\widehat{\boldsymbol f}]_\times\delta\boldsymbol\theta$,
subtracting nominal acceleration from true acceleration and retaining
first-order errors gives
\begin{align}
 \delta\dot{\boldsymbol p}&=\delta\boldsymbol v,\\
 \delta\dot{\boldsymbol v}&=-\widehat R_{wb}
       [\widehat{\boldsymbol f}]_\times\delta\boldsymbol\theta
       -\widehat R_{wb}\delta\boldsymbol b_a-\widehat R_{wb}\boldsymbol n_a,\\
 \delta\dot{\boldsymbol\theta}&=-[\widehat{\boldsymbol\omega}]_\times
       \delta\boldsymbol\theta-\delta\boldsymbol b_g-\boldsymbol n_g,\\
 \delta\dot{\boldsymbol b}_a&=\boldsymbol n_{ba},
 \qquad\delta\dot{\boldsymbol b}_g=\boldsymbol n_{bg}.
 \label{eq:eskf-dynamics}
\end{align}
For the attitude equation, differentiating
$R=\widehat R(I+[\delta\boldsymbol\theta]_\times)$ and using
$\dot R=R[\boldsymbol\omega]_\times$ leaves the commutator
$[\delta\boldsymbol\theta]_\times[\widehat{\boldsymbol\omega}]_\times
-[\widehat{\boldsymbol\omega}]_\times[\delta\boldsymbol\theta]_\times
=[\delta\boldsymbol\theta\times\widehat{\boldsymbol\omega}]_\times$,
which proves its sign.

Stack the errors as $(\delta p,\delta v,\delta\theta,\delta b_a,
\delta b_g)$ to obtain $\dot{\delta x}=F\delta x+G n$ from
\eqref{eq:eskf-dynamics}. For continuous noise spectral covariance $Q_c$,
\begin{equation}
 \dot P=FP+PF^{\mathsf T}+GQ_cG^{\mathsf T}.
\end{equation}
For frozen $F,G,Q_c$ over $\Delta t$, the exact discretization is
\begin{equation}
 \Phi=e^{F\Delta t},\qquad
 Q_d=\int_0^{\Delta t}e^{Fs}GQ_cG^{\mathsf T}e^{F^{\mathsf T}s}\,ds,
 \qquad P^-_{k+1}=\Phi P^+_k\Phi^{\mathsf T}+Q_d.
\end{equation}
For example, if
\begin{equation}
 \exp\left(\Delta t\begin{bmatrix}F&GQ_cG^{\mathsf T}\\0&-F^{\mathsf T}
                         \end{bmatrix}\right)
 =\begin{bmatrix}\Phi&B_d\\0&\Phi^{-\mathsf T}\end{bmatrix},
\end{equation}
then $Q_d=B_d\Phi^{\mathsf T}$, as follows from the block exponential
integral. $Q_c$ denotes a noise density, not a sampled sensor variance;
confusing these introduces an erroneous factor of sample interval.
The frozen-coefficient approximation must be evaluated at the actual
motion and sampling interval. No magnetometer or accelerometer gravity
update is implied by these equations.

For the scalar-first quaternion convention above, the rotation
kinematics also give
\begin{equation}
 \dot{\widehat q}=\tfrac12\widehat q\otimes
                  [0,\widehat{\boldsymbol\omega}],\qquad
 \widehat q_{k+1}=\widehat q_k\otimes
 \left[\cos\frac{\|\widehat\omega\|\Delta t}{2},
 \frac{\widehat\omega}{\|\widehat\omega\|}
                    \sin\frac{\|\widehat\omega\|\Delta t}{2}\right]
 \label{eq:quaternion-propagation}
\end{equation}
when the body angular rate is constant over the interval; take the
continuous zero-rate limit for the vector component. To derive the
differential equation, compose the attitude with the infinitesimal
right increment $[1,\widehat\omega\,dt/2]+O(dt^2)$ and divide the
quaternion difference by $dt$. For constant rate, exponentiating this
linear differential equation gives the trigonometric increment in
\eqref{eq:quaternion-propagation}, equivalently
$R_{k+1}=R_k\exp([\widehat\omega]_\times\Delta t)$. For a changing
rate, a single exponential of the rate integral is generally not exact,
because rotations about different axes do not commute. Use adequately
resolved integration or an appropriate coning correction, and verify
the integration error against the expected motion. Quaternion norm
renormalization corrects roundoff but does not correct integration bias.

\subsubsection{Landmark reprojection Jacobian and its sign}

At the nominal state define
$\boldsymbol p_b=\widehat R_{wb}^{\mathsf T}(\boldsymbol m-
\widehat{\boldsymbol p})$ and
$\boldsymbol p_c=R_{cb}(\boldsymbol p_b-\boldsymbol\ell_b)
=[X,Y,Z]^{\mathsf T}$. Pinhole differentiation yields
\begin{equation}
 J_\pi=\begin{bmatrix}f_x/Z&0&-f_xX/Z^2\\
                        0&f_y/Z&-f_yY/Z^2\end{bmatrix}.
\end{equation}
From \eqref{eq:right-attitude-error},
\begin{align}
 R_{wb}^{\mathsf T}(\boldsymbol m-\boldsymbol p)
 &\simeq(I-[\delta\boldsymbol\theta]_\times)
          (\boldsymbol p_b-\widehat R_{wb}^{\mathsf T}\delta\boldsymbol p)\\
 &=\boldsymbol p_b-\widehat R_{wb}^{\mathsf T}\delta\boldsymbol p
                  +[\boldsymbol p_b]_\times\delta\boldsymbol\theta.
\end{align}
Consequently
\begin{equation}
 H_p=-J_\pi R_{cb}\widehat R_{wb}^{\mathsf T},\qquad
 H_\theta=\boxed{+J_\pi R_{cb}[\boldsymbol p_b]_\times}.
 \label{eq:reprojection-jacobian}
\end{equation}
The positive attitude sign in Appendix H.2 is correct for this local
error convention. Reversing it without changing the convention would
be an error. With a lever arm, the cross-product uses $\boldsymbol p_b$,
not $\boldsymbol p_b-\boldsymbol\ell_b$, since the lever arm is fixed in
the body frame. Extrinsic uncertainty needs additional Jacobian columns.

For a known map point perturbation $\delta\boldsymbol m$,
$H_m=J_\pi R_{cb}\widehat R_{wb}^{\mathsf T}$. Known independent map
uncertainty can be propagated as $H_m\Sigma_m H_m^{\mathsf T}$ into the
measurement covariance. Shared datum errors and inter-landmark map
correlations require joint covariance or nuisance states.

For a height-above-terrain observation
$h(\boldsymbol p)=p_z-z_g(p_x,p_y)$, the position Jacobian is
$[-\partial_xz_g,-\partial_yz_g,1]$. A barometric altitude relative to
launch pressure instead measures altitude relative to that datum;
conversion to above-ground height must account for ground elevation
and datum uncertainty. Height bias is not removed by naming it AGL.

\subsubsection{Update, injection, and covariance reset}

Stack compatible measurements into an innovation $\nu=z-h(\widehat x)$
with covariance $V$ conditional on the state. With linearized Jacobian
$H$, independent measurement noise gives
\begin{align}
 S&=HP^-H^{\mathsf T}+V,&K_f&=P^-H^{\mathsf T}S^{-1},\\
 \widehat{\delta x}&=K_f\nu,&
 P_c&=(I-K_fH)P^-(I-K_fH)^{\mathsf T}+K_fVK_f^{\mathsf T}.
\end{align}
The last expression is the Joseph covariance update. For shared noise
or dependence between the prediction and measurement errors, cross
covariances must be included; the displayed equations assume they vanish.
Inject the Euclidean increments by addition and attitude by
\begin{equation}
 \widehat R^+=\widehat R^-\operatorname{Exp}
                     ([\widehat{\delta\boldsymbol\theta}]_\times).
\end{equation}
The error coordinates have now changed, so retaining $P_c$ unchanged is
not generally correct. If $a=\widehat{\delta\boldsymbol\theta}$, the new
attitude error is
\begin{equation}
 \delta\theta^+=\operatorname{Log}\{\operatorname{Exp}(-[a]_\times)
                                \operatorname{Exp}([\delta\theta^-]_\times)\}.
\end{equation}
Here $\operatorname{Log}$ returns the three-vector rotation increment
on its local principal branch, so
$\operatorname{Log}(\operatorname{Exp}([v]_\times))=v$ for
$\|v\|<\pi$. Differentiating at $\delta\theta^-=a$ gives the right Jacobian
\begin{align}
 J_r(a)&=\int_0^1\exp(-s[a]_\times)\,ds\\
 &=I-\frac{1-\cos\|a\|}{\|a\|^2}[a]_\times
       +\frac{\|a\|-\sin\|a\|}{\|a\|^3}[a]_\times^2
   =I-\tfrac12[a]_\times+O(\|a\|^2).
\end{align}
One way to see the derivative is
$\operatorname{Exp}([a+\varepsilon]_\times)
=\operatorname{Exp}([a]_\times)\operatorname{Exp}([J_r(a)\varepsilon]_\times)
+O(\|\varepsilon\|^2)$; left multiplication by
$\operatorname{Exp}(-[a]_\times)$ and the local logarithm then return
$J_r(a)\varepsilon$. Integrating Rodrigues' formula proves the closed
form. With true-minus-nominal Euclidean errors the reset matrix is
\begin{equation}
 T=\operatorname{diag}(I_3,I_3,J_r(a),I_3,I_3),\qquad
 P^+=T P_cT^{\mathsf T}.
\end{equation}
This also transforms attitude cross-covariances. Set the mean error
state back to zero after injection. Large residual rotations require
relinearization or iteration rather than trusting a first-order update.

\subsubsection{Visual odometry and its sign}

After rotational compensation, a nadir camera above stationary flat
terrain obeys $u=f(m_x-p_x)/h+c_x$. At constant height,
\begin{equation}
 \Delta p_x=-\frac h f\Delta u,\qquad
 v_x\simeq-\frac h f\frac{\Delta u}{\Delta t}.
 \label{eq:vo-sign}
\end{equation}
The minus sign applies when $\Delta u$ means the displacement of a
ground feature in image coordinates. A correlation shift defined as
the compensating image translation has the opposite sign, so code
must document which convention it uses. Appendix H.7's unsigned
formula is incomplete until its shift convention is specified.

With changing heights the exact nadir-plane relation is
\begin{equation}
 \Delta p_x=\frac{h_1}{f}(u_1-c_x)-\frac{h_2}{f}(u_2-c_x).
\end{equation}
Thus descent produces scale change even without horizontal motion.
Simple translation-only correlation of unscaled frames will confound
this expansion with drift. Use plane rectification at each height or
a scale-aware model; propagate height and attitude uncertainty.
Consecutive displacements share images and therefore have correlated
errors. Their velocity estimates are not independent new measurements
of the same filter that produced the rectification unless that
dependence is represented. Without fresh absolute map fixes, visual
odometry drift and datum uncertainty accumulate.

The temporal covariance can be derived directly. Suppose a tracked
feature has independent zero-mean image-location errors $\eta_k$ with
covariance $C_k$, and the known-height, known-rotation transformation to
horizontal ground position is $B_k$ (for the aligned nadir example,
$B_k=(h_k/f)I_2$). Its displacement error between images $k$ and $k+1$
is $\epsilon_k=B_k\eta_k-B_{k+1}\eta_{k+1}$. Therefore
\begin{align}
 \operatorname{Cov}(\epsilon_k)
 &=B_kC_kB_k^{\mathsf T}+B_{k+1}C_{k+1}B_{k+1}^{\mathsf T},\\
 \operatorname{Cov}(\epsilon_k,\epsilon_{k+1})
 &=-B_{k+1}C_{k+1}B_{k+1}^{\mathsf T}.
 \label{eq:vo-temporal-covariance}
\end{align}
Both equalities follow by expanding products and discarding independent
cross terms. The negative shared-image covariance is present even in
this idealized case. Velocity covariance divides the first expression
by $\Delta t_k^2$ and the adjacent cross-covariance by
$\Delta t_k\Delta t_{k+1}$. Tracking errors across frames, multiple
features, and pose-dependent rectification can add other cross terms;
thus the independent-location derivation is a baseline, not the full
pipeline model. Height error introduces additional derivatives of
$B_k$ and must be propagated as well.

If the prior state error $e$ and measurement noise $v$ are correlated,
let $C=\mathbb E[e v^{\mathsf T}]$ and write $\nu=He+v$. Expansion gives
\begin{align}
 S&=HPH^{\mathsf T}+HC+C^{\mathsf T}H^{\mathsf T}+V,\\
 \operatorname{Cov}(e,\nu)&=PH^{\mathsf T}+C,\\
 K_f&=(PH^{\mathsf T}+C)S^{-1},\\
 P_c&=P-(PH^{\mathsf T}+C)S^{-1}(HP+C^{\mathsf T}).
 \label{eq:correlated-measurement-update}
\end{align}
The gain follows by minimizing the covariance of $e-K_f\nu$, or by
Gaussian conditioning. These equations show precisely which omitted
terms can cause double use of old visual information. For temporal
measurement correlation, whiten the complete noise model, augment the
state with its correlated components, or use a joint update with the
block covariance; treating every velocity increment as independent
does not reproduce that update.

\subsection{Error budget, information bounds, and filter consistency}

\subsubsection{A design sensitivity calculation is not posterior covariance}

The proposal's scalar budget assumes independent source errors and a
specified one-dimensional projection. Using $r_m$ for map metres/pixel
to avoid confusing it with gravity, it is
\begin{equation}
 \sigma_p^2\simeq(h\sigma_\theta)^2+
       \frac{r_m^2\sigma_c^2}{n}+\sigma_{\mathrm{datum}}^2+
       (h\omega\sigma_\tau)^2+(\sigma_z\tan\phi)^2.
 \label{eq:scalar-budget}
\end{equation}
This follows from a linear error $e=\sum_j a_j\eta_j$ and
$\operatorname{Var}(e)=\sum_j a_j^2\operatorname{Var}(\eta_j)$ only when
the source covariances vanish. In general, with $e=A\eta$,
\begin{equation}
 \Sigma_e=A\Sigma_\eta A^{\mathsf T}
          =\sum_{j,k}A_j\operatorname{Cov}(\eta_j,\eta_k)A_k^{\mathsf T}.
\end{equation}
Matching noise averages down as $1/n$ only for independent, equal-quality,
translation-only observations. Repeated texture, image warp, timing,
terrain geometry, and map datum are often common across landmarks.
For equal correlation $\rho$ between scalar matching errors, the mean
variance is $\sigma^2\{(1-\rho)/n+\rho\}$, obtained by summing the
$n$ diagonal and $n(n-1)$ off-diagonal covariance entries.

For $r_m=0.5\,\mathrm{m/pixel}$, $\sigma_c=0.2$ pixels and $n=50$,
the independent translation matching term is
$0.0141421\,\mathrm m$. The assumption is much more consequential than
its rounded value. A common map offset remains in every fix and does
not shrink by adding landmarks or frames.

The ESKF already estimates position and attitude jointly. Its posterior
$P_{pp}$ contains the effects of marginalized attitude uncertainty and
correlation with position. Adding $h^2P_{\theta\theta}$ to that same
$P_{pp}$ generally double-counts uncertainty. Either use the joint
posterior and propagate it through the desired output, or construct a
separate, explicitly defined design estimator and its source sensitivities.
Map datum and calibration nuisance uncertainties not in the filter
must be added with their correct cross-covariance or included as states.

The nominal attitude--datum crossover is obtained by setting
$h\sigma_\theta=\sigma_{\mathrm{datum}}$:
\begin{equation}
 h^*=\frac{\sigma_{\mathrm{datum}}}{\sigma_\theta}
      =\frac{2.5}{0.5\pi/180}=286.479\,\mathrm m.
\end{equation}
The units require radians. Inferring attitude precision from an observed
crossover is not unique unless the datum, terrain, timing, correlations,
and altitude-dependent geometry are separately controlled. If tilt is
estimated from fixed angular geometry with fixed pixel noise, its
posterior precision may itself vary with height and feature support.

In \eqref{eq:scalar-budget}, taking the terrain contribution to be exactly
$1\,\mathrm m$, timing $0.5\,\mathrm{ms}$ at $40^\circ/\mathrm s$, and
the matching value above gives $2.73470$ to $6.11774\,\mathrm m$ at
$300\,\mathrm m$ for $(\sigma_\theta,\sigma_{\mathrm{datum}})$ from
$(0.3^\circ,2\,\mathrm m)$ to $(1^\circ,3\,\mathrm m)$. At
$100\,\mathrm m$ the corresponding endpoints are $2.29686$ and
$3.61215\,\mathrm m$. These are worked sensitivity estimates, not
predictions validated on hardware. State whether an error is a
per-axis standard deviation, horizontal RMS $\sqrt{\operatorname{tr}
\Sigma_p}$, radial quantile, or bias bound: they are different quantities.

\subsubsection{Fisher information and nuisance marginalization}

For fixed correct correspondences with independent Gaussian observations
$z_i=h_i(x)+\varepsilon_i$, covariance $V_i$ independent of $x$, the
log-likelihood is a sum of weighted squared residuals. Taking the
expected outer product of the score gives
\begin{equation}
 J_{\mathrm{lik}}=\sum_i H_i^{\mathsf T}V_i^{-1}H_i.
\end{equation}
For a differentiable identifiable deterministic parameter, the
Cram\'er--Rao inequality follows by differentiating estimator
unbiasedness and applying the covariance Schur complement to the
estimator error and likelihood score:
\begin{equation}
 \operatorname{Cov}(\widehat x)\succeq J_{\mathrm{lik}}^{-1}.
\end{equation}
The regularity assumptions allow differentiation under the integral,
and the unbiasedness derivative must be identity. A biased estimator,
selected associations, wrong measurement model, or singular information
invalidates this simple statement.

With a Gaussian linear prior, $J_0=P_0^{-1}$ and linear measurements,
the posterior precision is exactly $J_0+J_{\mathrm{lik}}$. In a nonlinear
EKF it is a local linearized posterior approximation. Adding a prior
does not turn this into a frequentist bound for \emph{any} unbiased
estimator without specifying a Bayesian bound and its assumptions.
For correlated observations, stack $H$ and use $J=H^{\mathsf T}V^{-1}H$;
the independent sum is not valid unless conditional independence holds.

Partition information between parameters of interest $x$ and nuisance
$b$. If $J_{bb}$ is invertible, the effective information is
\begin{equation}
 J_{x,\mathrm{eff}}=J_{xx}-J_{xb}J_{bb}^{-1}J_{bx}.
\end{equation}
Block Gaussian elimination proves this formula and
$(J^{-1})_{xx}=J_{x,\mathrm{eff}}^{-1}$. It explicitly shows why
pretending tilt, datum, or calibration is known overstates accuracy.
If a gauge direction is unobservable, restrict to an identifiable
subspace or add a justified independent datum measurement. A
pseudoinverse alone must not be reported as a finite error bound in an
unobservable direction.

\subsubsection{Innovation gates and consistency against imperfect truth}

For an unselected correct landmark with Gaussian innovation
$\nu\sim N(0,S)$ and rank two covariance, whitening gives
\begin{equation}
 \nu^{\mathsf T}S^{-1}\nu\sim\chi_2^2,\qquad
 \chi_2^2(0.997)=-2\log(0.003)=11.618286.
\end{equation}
The equality uses the radial squared length of two independent unit
normal variables: in polar coordinates its density is
$\tfrac12e^{-q/2}$ for $q\geq0$. A threshold rejects an inlier with
probability $0.003$ under this model. It does not bound acceptance of
a bad association. Searching many candidate peaks, consensus selection,
and prior dependence can alter the distribution. Validate the complete
selection and gating pipeline, including rejected cases. Several
landmarks with shared noise should be gated jointly or treated with
their joint covariance.

For true state error $e\sim N(0,P)$ in an $n$-dimensional observable
state, $\epsilon=e^{\mathsf T}P^{-1}e\sim\chi_n^2$. For $N$ independent
realizations at a comparable epoch,
\begin{equation}
 \sum_{r=1}^N\epsilon_r=N\overline\epsilon\sim\chi_{nN}^2.
\end{equation}
A two-sided significance level $\alpha$ gives average bounds
$[\chi_{nN}^2(\alpha/2)/N,
\chi_{nN}^2(1-\alpha/2)/N]$. High values indicate excess error relative
to stated covariance; low values can indicate conservative covariance,
dependence, or selection effects. Successive flight epochs are correlated
and are not $N$ independent realizations. One flight cannot be turned
into many independent trials merely by increasing sample rate.

GPS position cannot supply the true full 15-dimensional ESKF state.
For the horizontal difference $d=\widehat p_{\mathrm{TRN}}-
\widehat p_{\mathrm{GPS}}$, use the observed dimension and
\begin{equation}
 \Sigma_d=\Sigma_{\mathrm{TRN}}+\Sigma_{\mathrm{GPS}}
               -C_{TG}-C_{TG}^{\mathsf T},
\end{equation}
where $C_{TG}$ is the cross-covariance of the two estimation errors.
Even with GPS excluded from navigation, datum survey, time alignment,
or calibration can make errors correlated. Include map datum uncertainty
if it is not already in $\Sigma_{\mathrm{TRN}}$, and distinguish a measured
bias from random error. Then $d^{\mathsf T}\Sigma_d^{-1}d$ tests
consistency of the difference model; it is not pure navigation NEES
unless the truth is sufficiently accurate and independent.

\subsection{Horizontal protection level and integrity limitations}

Assume horizontal error $e\sim N(0,\Sigma_p)$ with
$\Sigma_p\succ0$. Write $e=\Sigma_p^{1/2}\xi$,
$\xi\sim N(0,I_2)$, and let
$q_\alpha=\chi_2^2(1-\alpha)=-2\log\alpha$.
The ellipsoid
\begin{equation}
 \mathcal E_\alpha=\{e:e^{\mathsf T}\Sigma_p^{-1}e\leq q_\alpha\}
\end{equation}
has probability $1-\alpha$. Since
$\|e\|^2\leq\lambda_{\max}(\Sigma_p)\|\xi\|^2$, the circle
\begin{equation}
 \mathrm{PL}=\sqrt{-2\log\alpha}\,
                          \sqrt{\lambda_{\max}(\Sigma_p)}
 \label{eq:PL-proof}
\end{equation}
contains that ellipsoid and has exceedance probability at most
$\alpha$. The circle probability is exactly $1-\alpha$ for isotropic
covariance; for anisotropic covariance it can be conservative.
At $95\%$ containment $k=\sqrt{-2\log0.05}=2.44774683$.
Using $k=3$ in isotropic two-dimensional error gives tail
$e^{-9/2}=0.011108997$, not the one-dimensional ``three sigma'' tail.

If an unknown deterministic bias has a justified bound $\|b\|\leq B$,
the triangle inequality gives $\mathrm{PL}=B+
\sqrt{q_\alpha\lambda_{\max}}$. If there is an unmodelled failure
probability $p_f$, a conditional Gaussian bound provides, at best,
total exceedance bound $p_f+\alpha$ by a union bound. A covariance
circle alone supplies no protection against unbounded datum error,
systematic model error, or wrong matches.

Empirical exceedance counts require uncertainty on their estimated
rate. With $M$ independent trials and zero exceedances, the one-sided
$1-\beta$ binomial upper bound is
\begin{equation}
 p_{\mathrm{exceed}}\leq1-\beta^{1/M},
\end{equation}
because the probability of zero under rate $p$ is $(1-p)^M$. Correlated
epochs do not meet this independence assumption. Comparing against GPS
also measures the error of the difference, not exact true TRN error.
Report integrity as a model and an experimental diagnostic; the
proposal does not establish an aviation certification claim.

\subsection{Map gauge, robust registration, and certification}

\subsubsection{The map datum gauge}

If all map coordinates and the camera position are shifted by the same
$d$, then $(m_i+d)-(C+d)=m_i-C$, so every predicted pixel is unchanged.
More generally, changing map/world coordinates by a rigid transform
$m_i'=Qm_i+d$, $C'=QC+d$, and $R_{cw}'=R_{cw}Q^{\mathsf T}$ leaves
the projection unchanged. Gravity and other references can constrain
some frame rotations, but imagery alone cannot determine an unknown
absolute map origin. An external surveyed datum is required.
Represent its uncertainty as a shared quantity across images and
landmarks. A fixed known map removes the gauge algebraically; its
georeferencing accuracy remains a separate physical uncertainty.

\subsubsection{Exact truncated-least-squares reformulation}

Let $q=[a,b,t_x,t_y]^{\mathsf T}$ represent complex scale-rotation
$z=a+ib$ and translation. For point $u_i=(u_{ix},u_{iy})$, define
\begin{equation}
 A_i=\begin{bmatrix}u_{ix}&-u_{iy}&1&0\\
                       u_{iy}&u_{ix}&0&1\end{bmatrix},\qquad
 r_i=m_i-A_iq.
\end{equation}
The similarity residual is affine, but the robust objective
\begin{equation}
 F(q)=\sum_i\min\left(\frac{\|r_i\|^2}{\sigma_i^2},c_i^2\right)
\end{equation}
is nonconvex. Exactly,
\begin{equation}
 F(q)=\min_{w_i\in\{0,1\}}\sum_i
       \left[w_i\frac{\|r_i\|^2}{\sigma_i^2}+(1-w_i)c_i^2\right].
\end{equation}
For fixed $q$ each independent switch chooses the smaller term. However,
$w_i\|m_i-A_iq\|^2$ has total polynomial degree three in $(w_i,q)$.
Thus the proposal's statement that binary switches immediately give a
quadratic programme is incomplete.

An explicit exact quadratic lifting introduces $v_i=w_iq$ and constraints
\begin{equation}
 w_i^2-w_i=0,\qquad v_{ij}-w_iq_j=0\quad(j=1,\ldots,4).
\end{equation}
Then solve the quadratically constrained quadratic programme
\begin{equation}
 \min_{q,w,v}\sum_i\left[
  \frac{\|w_im_i-A_iv_i\|^2}{\sigma_i^2}+(1-w_i)c_i^2\right].
 \label{eq:lifted-TLS}
\end{equation}
When $w_i=0$ the constraint forces $v_i=0$, giving $c_i^2$.
When $w_i=1$, $v_i=q$ and the cost is the original squared residual.
These two cases prove equivalence in both directions. Additional
physically justified scale constraints may be imposed through
$s_{\min}^2\leq a^2+b^2\leq s_{\max}^2$; translation/search bounds
must also be stated if used. Without a nonzero scale restriction,
$z=0$ is allowed algebraically although it is a collapsed transformation,
not an invertible similarity.

\subsubsection{What an SDP certificate requires}

Collect the variables and a homogeneous leading $1$ into $x$.
Every objective and quadratic constraint can be written as
$x^{\mathsf T}Qx$ and $x^{\mathsf T}A_jx=b_j$, including linear terms
through the leading entry. Set $X=xx^{\mathsf T}$, enforce
$X_{00}=1$, and replace quadratic forms by traces. Dropping
$\operatorname{rank}X=1$ gives an SDP with $X\succeq0$, linear
trace constraints, and linear objective. Any exact original feasible
point generates an SDP feasible matrix with the same objective, so
\begin{equation}
 L_{\mathrm{SDP}}\leq F_{\mathrm{global}}.
\end{equation}
If a globally optimal SDP solution has rank one and its factor with
$x_0=1$ satisfies every original constraint, it provides an original
feasible point with cost $L_{\mathrm{SDP}}$. Therefore
\begin{equation}
 L_{\mathrm{SDP}}\leq F_{\mathrm{global}}\leq F(x)=L_{\mathrm{SDP}},
\end{equation}
which proves global optimality. This is the mathematical content of a
lossless relaxation certificate. The primary work by Yang and Carlone
develops semidefinite approaches to robust geometric estimation.\footnote{
Heng Yang and Luca Carlone, \emph{Certifiably Optimal Outlier-Robust
Geometric Perception}: \url{https://arxiv.org/abs/2109.03349}.}

Rank one \emph{by itself} is insufficient: a suboptimal rank-one feasible
point is just a candidate. A floating-point eigenvalue ratio is a
numerical diagnostic, not a proof of exact rank, feasibility, or
optimality. A practical claim must report primal constraint residuals,
original-variable feasibility, recovered original objective, and a
valid dual lower bound. For an equality-constrained SDP, a dual vector
$y$ satisfying $Q-\sum_j y_jA_j\succeq0$ gives the lower bound
$\sum_jy_jb_j$. Together with a feasible candidate of cost $U$, it
certifies suboptimality at most $U-L$. Inequality constraints require
the corresponding sign restrictions on multipliers. Tolerance-based
claims must be named numerical certificates; exact certificates require
validated arithmetic or exact bounds. A relaxation need not be tight,
and a rank-one optimum need not be returned even when one exists.
Finally, global optimality of the stated robust cost does not prove
that the map correspondence is physically correct.

\subsection{Landmark selection: normalized objective and greedy proof}

Assume a fixed observable parameterization, a positive definite baseline
information matrix $J_0$, and fixed positive semidefinite per-landmark
information $J_i$. Define
\begin{equation}
 J_S=J_0+\sum_{i\in S}J_i,\qquad
 f(S)=\log\det J_S-\log\det J_0.
 \label{eq:normalized-information}
\end{equation}
Then $f(\varnothing)=0$. The subtraction is essential for the
approximation guarantee and removes the arbitrary log-determinant
constant arising from coordinate units.

For $i\notin S$, the determinant identity gives
\begin{equation}
 \Delta_i(S)=f(S\cup\{i\})-f(S)
 =\log\det\{I+J_i^{1/2}J_S^{-1}J_i^{1/2}\}\geq0.
\end{equation}
The equality follows from $\det(I+AB)=\det(I+BA)$ and factorization by
$J_S$. If $A\subseteq B$, then $J_A\preceq J_B$ and inversion reverses
the positive definite order, $J_A^{-1}\succeq J_B^{-1}$.
Congruence with $J_i^{1/2}$ preserves that order. The log determinant
is increasing on positive definite matrices: along a PSD increment
$D$, its derivative is $\operatorname{tr}[(M+tD)^{-1}D]\geq0$.
Therefore
\begin{equation}
 \Delta_i(A)\geq\Delta_i(B),
\end{equation}
proving diminishing returns, hence submodularity and monotonicity.

For a cardinality constraint $|S|\leq k$, let $O$ be optimal and
$S_t$ the greedy set after $t$ additions. By monotonicity and repeated
diminishing returns,
\begin{align}
 f(O)-f(S_t)&\leq f(O\cup S_t)-f(S_t)\\
 &\leq\sum_{i\in O\setminus S_t}\Delta_i(S_t)
 \leq k\{f(S_{t+1})-f(S_t)\}.
\end{align}
Thus the remaining gap contracts by at least $1-1/k$ at each addition:
\begin{equation}
 f(S_k)\geq[1-(1-1/k)^k]f(O)\geq(1-1/e)f(O).
\end{equation}
This is the normalized gain guarantee, not a guarantee that the raw
$\log\det J_{S_k}$ is at least $(1-1/e)$ times the raw optimum.
The original approximation theorem is due to Nemhauser, Wolsey, and
Fisher.\footnote{\emph{An analysis of approximations for maximizing
submodular set functions--I} (1978):
\url{https://thibaut.horel.org/submodularity/papers/nemhauser1978.pdf}.}

Independent observations conditional on an augmented state can supply
fixed additive $J_i$. Correlated measurement noise, selection-dependent
linearization, or optimizing a pose-only determinant after marginalizing
shared nuisance states need not meet that additive model; do not assert
this proof then applies without a separate argument. A singular $J_0$
makes the normalized log determinant undefined. Remove gauge directions
or use a justified prior; a fictitious prior conceals an unobservable
mode. Unequal landmark costs require a different budget-constrained
algorithm and guarantee. Greedy can prefer outer features in the simple
tilt model, but maximizing full pose information can also select inner
features, directions that fill a weak mode, or lower-noise landmarks.
``Always chooses edge landmarks'' is not a theorem.

\subsection{Uncertainty-aware notional safe-target selection}

Let $c(p)$ be a bounded nonnegative map cost and let terminal displacement
error have distribution $e\sim N(0,\Sigma_{\mathrm{land}})$ under the
specified predictor. If $p$ is the intended mean target, expected cost is
\begin{equation}
 c_\Sigma(p)=\mathbb E[c(p+e)]
 =\int_{\mathbb R^2}c(p+e)\,
                      \mathcal N(e;0,\Sigma_{\mathrm{land}})\,de.
 \label{eq:expected-target-cost}
\end{equation}
For the indicator of a hazard set $\mathcal H$, this equals the modeled
probability of landing in a hazard. It is a convolution, not merely
an unspecified inflation radius. Use terminal covariance, map uncertainty,
and actuator/model uncertainty appropriate to the target decision;
current position covariance alone need not equal terminal uncertainty.

A robust alternative uses the confidence error ellipsoid
$\mathcal E_\alpha$ defined above. Require
\begin{equation}
 (p+\mathcal E_\alpha)\cap\mathcal H=\varnothing
 \quad\Longleftrightarrow\quad
 p\notin\mathcal H\oplus(-\mathcal E_\alpha),
\end{equation}
where $\oplus$ is Minkowski addition. Then all errors in the ellipsoid
avoid hazards, yielding modeled hazard probability at most $\alpha$.
For a centered symmetric ellipse $-\mathcal E_\alpha=\mathcal E_\alpha$.
Add a bounded datum/hazard-boundary uncertainty set by another Minkowski
sum. Gaussian support is unbounded, so finite inflation is a probability
statement, not zero-risk clearance. Map completeness remains an
unverified assumption until inspected.

For a notional controllable canopy, let available horizontal
air-relative control speed be $u_{\max}(t)$. The triangle inequality
gives a displacement bound
\begin{equation}
 \|\Delta p_{\mathrm{control}}\|
 \leq\int_t^{t_{\mathrm{impact}}}u_{\max}(t')\,dt'.
\end{equation}
If $u_{\max}=\kappa v_d$, vertical descent is monotone with
$dh/dt=-v_d$, and control can point freely in the plane without delay,
then the bound is $R_{\mathrm{div}}=\kappa h$. Under those idealized
conditions every point in the disc can be generated by suitable bounded
control. At $\kappa=0.3$ this gives $90\,\mathrm m$ at $300\,\mathrm m$
and $30\,\mathrm m$ at $100\,\mathrm m$. If turn-rate limits,
deployment transients, directional constraints, or spatially varying
wind affect the control path, this disc can overstate reachability.

Define the compact, nonempty modeled feasible set $\mathcal D$ and choose
\begin{equation}
 p^*\in\operatorname*{arg\,min}_{p\in\mathcal D}c_\Sigma(p).
\end{equation}
A continuous cost on a compact set attains a minimum. On a finite
raster, search all admissible cells and specify a deterministic tie
rule; quantify discretization and boundary handling. If no acceptable
cell exists, report infeasibility rather than outputting a falsely
safe target. The actual passive CanSat has no commanded divert, so
$\kappa=0.3$ describes a hypothetical decision experiment, not its
physical reachable set. Agreement between TRN- and GPS-driven decisions
tests sensitivity to navigation, not proof that an unflown control
would reach the target.

\subsection{Executed numerical checks and remaining validation}

The following checks were executed with NumPy and Python, using
$\texttt{default\_rng}(20260930)$ for random inputs. No SciPy, SDP solver,
flight imagery, or flight measurements were used. The random rotations
were generated with Rodrigues' formula. The reproducible module
\texttt{sim/maths/test\_navigation\_extra.py} prints the numerical report
and executes thirteen unit tests; its derivatives, rotations, and FFT
checks also use the adjacent mathematical reference models. Reprojection tests used
$K=(f_x=f_y=640,c_x=640,c_y=400)$, nominal camera height $300\,\mathrm m$,
100 random ground points and attitude perturbations, a nonidentity
camera extrinsic rotation, and body lever arm $(0.04,-0.03,0.015)\,\mathrm m$.
Symmetric finite differences used step $10^{-5}$ for projection, with
rotations perturbed on the right. These are algebraic consistency checks.

\begin{center}
\begin{tabular}{@{}p{0.55\linewidth}p{0.36\linewidth}@{}}
\toprule
Executed check & Result\\
\midrule
Homography versus direct projection, maximum pixel difference
 & $2.03\times10^{-13}$ pixels\\
Position reprojection derivative, maximum component error
 & $1.02\times10^{-8}$ pixels/m\\
Local attitude reprojection derivative, maximum component error
 & $6.96\times10^{-8}$ pixels/rad\\
Exact reset derivative at $a=(0.3,-0.2,0.1)$ rad,
step $10^{-6}$, versus $J_r(a)$
 & $3.07\times10^{-11}$ maximum component error\\
FFT NCC versus direct NCC, all 117 valid windows of a
$13\times16$ random image and $5\times4$ random template
 & $4.45\times10^{-16}$ maximum score difference\\
Gaussian radial $95\%$ circle, 200000 independent draws
 & Observed containment $0.9492$; specified probability $0.95$\\
\bottomrule
\end{tabular}
\end{center}

For the Gaussian circle experiment, the binomial standard deviation of
the sample proportion under the specified model is
$\sqrt{0.95(0.05)/200000}=0.0004873$. The observed difference $-0.0008$
is about $-1.64$ such standard deviations. This is compatible with
sampling variation and verifies the simulated distribution; it gives
no evidence for the Gaussianity of real matching errors.

For \eqref{eq:tilt-information}, with $f=640$ and
$\sigma_{\mathrm{px}}=0.2$, the normalized landmark sets
$(-0.02,-0.01,0,0.01,0.02)$,
$(-1,-1,1,1)$, and $(-1,-0.5,0,0.5,1)$ gave marginalized attitude
information $1.4336$, $0$, and $8.96\times10^6$ rad$^{-2}$,
respectively. These confirm strong near-centre ill-conditioning and
the equal-radius degeneracy in the one-dimensional model. They are not
full-camera attitude performance predictions.

A separate normalized-logdet check generated eight fixed PSD landmark
matrices $J_i=H_i^{\mathsf T}H_i$ from $2\times3$ normal matrices and
used $J_0=\operatorname{diag}(0.2,0.4,0.6)$. All diminishing-return
comparisons between a subset and a one-element larger subset were
checked; the maximum positive violation was zero in floating-point
arithmetic. Greedy selection of three landmarks gave normalized value
$8.573753638$, while exhaustive search over all $\binom83=56$ triples
gave $8.598882151$, a ratio $0.997077700$. This is a numerical example
of the proved guarantee, not a proof for every matrix set or a statement
about real image landmarks.

Additional passing tests compare constant-rate quaternion propagation
against an independent matrix exponential and Runge--Kutta integration,
the right Jacobian against numerical integration of
$\int_0^1\exp(-s[a]_\times)ds$, and the shared-image covariance
\eqref{eq:vo-temporal-covariance} against $300000$ independently generated
triples of image errors. They also check signed odometry with changing
height, undefined constant-template NCC, the two-dimensional quadratic
peak counterexample, integer RANSAC counts, the exact two binary cases
of the TLS lifting, and the Gaussian innovation gate. The latter
rejected $0.00314$ of the $200000$ synthetic Gaussian innovations at
the specified $0.003$ nominal rate. No test substitutes for the proof
of the corresponding general identity, and no test claims an SDP solve.

No semidefinite certificate has been computed here; the exact lifted
formulation and certificate conditions above specify what would need
to be checked. Before interpreting navigation simulations as physical
validation, run the same pipeline on independently rendered terrain
with elevation, calibration error, timing offset, motion blur, shared
map datum error, and repetitive textures. Retain failed associations
and rejected updates. Then validate with surveyed ground imagery and
independent attitude and position references. The required flight
accuracy, tilt observability under real swing, covariance consistency,
and match integrity remain hypotheses until such evidence exists.
